% ============================================================
%  ch06  第二讲 §2.2 漫谈、§2.3 对偶 Schubert「簇」与横截片
%  原文行 1467–1651
% ============================================================

\subsection{插叙：$\Gr$ 中的一些子归纳概形}\label{digression}

这里插叙一种构造 $\Gr$ 中若干重要子簇的方法。本小节在余下部分不再使用。

设 $\varphi: \Gr\to \Gr$ 是 $k$ 上态射。对 $\ga\in LG$ 与 $\mu$，定义 $\Gr$ 的闭子归纳概形
\[X(\mu,\ga\varphi)=\left\{ x\in \Gr\mid \inv(x,\ga\varphi(x))\leq \mu\right\}.\]
确切地说，写 $x=g L^+G\in \Gr$，则
\[X(\mu,\ga\varphi)= \left\{g L^+G\mid g^{-1}\ga\varphi(g) \in (LG)_{\leq \mu}\right\}/L^+G.\]

这里给出几个等价定义。

\begin{enumerate}
\item[(1)] $X(\ga\varphi,\mu)$ 由下述笛卡尔方块定义
\[\begin{CD}
X(\mu,\ga\varphi)@>>> \Gr\tilde\times\Gr_\mu\\
@VVV@VVV\\
\Gr@>1\times \ga\varphi>> \Gr\times \Gr.
\end{CD}\]

\item[(2)] 设 $\varphi$ 由保持 $L^+G$ 的态射 $\varphi:LG\to LG$ 诱导。考虑 $LG$ 被 $\varphi$ 扭变后的共轭作用，即 $\on{Int}(g)\cdot^\varphi g'= g^{-1}g' \varphi(g)$。考虑叠的态射
\[\pi: [(LG)_\mu/^\varphi L^+G]\to [LG/^\varphi LG].\]
则 $X(\mu,\ga\varphi)= \pi^{-1}(\ga)$。
\end{enumerate}

\begin{rmk}
许多时候，商 $[LG/^\varphi LG]$ 虽然性质不佳，仍应被看作某些重要代数结构的模空间。
\end{rmk}

下面是上述构造的几个具体例子。

\begin{enumerate}
\item[(i)] 取 $\varphi=\id$ 为恒等映射。则 $X(\ga,\mu)$ 是群版本仿射 Springer 纤维（见 \cite[\S~2.7.3]{Yu}）。此时 $[LG/LG]$ 分类 $G$ 的共轭类。该构造也有 Lie 代数版本，给出（通常的）仿射 Springer 纤维（见 \cite[\S~2.2]{Yu}）。

\item[(ii)] 设 $k=\bF_q$，$\varphi=\sigma$ 是 $\Gr$ 的 $q$-Frobenius。则 $X(\mu,\ga\varphi)$ 通常记作 $X(\mu,b)$，称为\textbf{仿射 Deligne--Lusztig 簇}（最早在 \cite{Ra} 中引入）。此时 $[LG/^{\sigma}LG]$ 分类 $G(\bar{k}((t)))$ 的 $\sigma$-共轭类，它们与 Drinfeld shtuka 密切相关。

\item[(iii)] 仍设 $k=\bF_p$。设 $\varphi:\Gr\to \Gr$ 由绝对 Frobenius $k[[t]]\mapsto k[[t]]$，$\sum a_it^i\mapsto \sum a_i^pt^{ip}$ 诱导。把 $\ga\varphi$ 记作 $\Phi$。则 $X(\Phi,\mu)$ 称为 \textbf{Kisin 簇}，而 $[LG/^{\sigma}LG]$ 分类 Galois 表示 $\Gal(\overline F/ F)\to G(\bF_p)$（见 \cite{PR2}）。

\item[(iv)] 设 $k=\bC$，$\varphi$ 由环路旋转 $t\mapsto qt$ 诱导，其中 $q\in \bG_m^{\on{rot}}$。则 $[LG/^{\varphi}LG]$ 与椭圆曲线 $E_q=\bC^\times/q^\bZ$ 上的 $G$-丛密切相关（Looijenga、Baranovsky 与 Ginzburg \cite{BG}）。

\item[(v)] 不取 $\varphi\in \bG_m^{\on{rot}}$，也可以让 $\varphi\in \Aut^{++}(D)$。该构造的 Lie 代数版本在 \cite{FZ} 中研究过。
\end{enumerate}

每种情形下都值得问几个基本问题：
\begin{enumerate}
\item 该簇何时有限维？
\item 若有限维，是否有维数公式？各不可约分支是否等维？
\item 如何参数化连通分支？
\item 如何参数化不可约分支？
\end{enumerate}

这些问题在情形 (i) 与 (ii) 中已了解得较好（已知结果的概述见 \cite{Gortz, Yu} 及其所引文献），其他情形则基本未涉及。

\begin{rmk}由于 \S\ref{p-adic} 的讨论，仿射 Springer 纤维与仿射 Deligne--Lusztig 簇现在对 $p$-进群也有定义。\end{rmk}

\subsection{对偶 Schubert「簇」与横截片}

对偶 Schubert「簇」与 \S\ref{II:Sch var} 引入的 Schubert 簇 $\Gr_{\leq \mu}$ 很不一样。事实上，它们并不是真正的簇，而且无穷维（但余维有限）。

先定义群预层
\[L^-G(R)=G(R[t^{-1}]).\]
\S\ref{Kac-Moody} 将说明 $LG$ 与 Kac--Moody 群密切相关；从这个观点看，$L^-G$ 是 $L^+G$ 的「对偶抛物子」。还有另一种看法：$L^-G$ 是 \eqref{Gout} 中将引入的群 $\underline G^{X^*}$ 的特例（对应 $X=\bP^1$、$x=0$）。特别地，由下文引理 \ref{rep of Gout}，群 $L^-G$ 可由归纳概形表示。

如 \S\ref{II:Sch var} 那样固定 $T\subset B\subset G$。则 $G(F)$ 有 \textbf{Birkhoff 分解}
\begin{equation}\label{Birk decom}
G(F)=\bigsqcup_{\mu\in \xcoch(T)^+} G(k[t^{-1}])t^\mu G(\mO),
\end{equation}
它诱导
\begin{equation}\label{Birk decom 2}
L^-G(k)\backslash G(F)/G(\mO)\cong \xcoch(T)^+.
\end{equation}
与 Cartan 分解 \eqref{Cartan decomp2} 类似，\eqref{Birk decom 2} 不依赖任何选取。一般 $G$ 情形的证明见 \cite[Lemma 4]{Fa}。

把 $L^-G$ 在 $\Gr$ 上的轨道称为\textbf{对偶 Schubert 胞腔}。由 Birkhoff 分解 \eqref{Birk decom 2}，它们形如
\[\Gr^\mu=L^-G\cdot t^\mu \subset \Gr,\]
并由 $\xcoch(T)^+$ 参数化。令
\[\Gr^{\geq \mu}=\bigsqcup_{\la\geq \mu}\Gr^\la,\]
称为\textbf{对偶 Schubert 簇}。下述命题说明这个术语是合理的。

\begin{prop}\label{geom of opp Sch}~
\begin{enumerate}
\item $\Gr^\mu$ 是 $\Gr$ 的局部闭子归纳概形。

\item $\Gr_\mu\cap \Gr^\la\neq \emptyset$ 当且仅当 $\mu\geq \la$。此外 $\Gr_\mu\cap \Gr^\mu= G\cdot t^\mu\cong G/P_\mu$，它在旋转环面 $\bG_m^{\on{rot}}$ 下不动。

\item $\Gr^{\geq \mu}$ Zariski 闭，且以 $\Gr^\mu$ 为稠密开子集。

\item $\Gr^{\geq \mu}$ 的余维为 $(2\rho,\mu)-\dim G/P_\mu$。
\end{enumerate}
\end{prop}

\begin{proof}
(1) $t^\mu$ 在 $L^-G$ 中的稳定化子是 $L^-G\cap t^\mu L^+G t^{-\mu}$，它是有限维子群。把 $L^-G$ 写成 $\varinjlim K_i$，使 $K_i\supset L^-G\cap t^\mu L^+G t^{-\mu}$，则每个 $K_it^\mu$ 在 $\Gr$ 中局部闭，故 $\Gr^\mu=\varinjlim K_it^\mu$ 是归纳概形。还要注意
\begin{equation}\label{dim opp Sch}
\dim(L^-G\cap t^\mu L^+G t^{-\mu})=(2\rho,\mu)+\dim P_\mu,
\end{equation}
这可由计算其李代数的维数得到，论证同命题 \ref{geom of Sch}。

(2) 留作练习。

(3) 先说明 $\Gr^0$ 在 $\Gr$ 的单位连通分支中开。注意有求值映射
\[L^-G\to G,\quad g\mapsto g\bmod t^{-1}.\]
以上述映射的核定义 $L^{<0}G=\ker(L^-G\to G)$。注意 $\Gr^0= L^{<0}G\cdot e$。

\begin{lem}\label{big cell}
乘法映射
\[L^{+}G\times L^{<0}G\to LG\]
是开嵌入。特别地，$L^{<0}G\cong L^{<0}G\cdot e\subset \Gr$ 是开浸入。
\end{lem}

\begin{proof}
注意 $L^{<0}G\cap L^+G=\{1\}$，且在李代数层面 $L\frakg=L^+\frakg\oplus L^{<0}\frakg$，其中 $L^+\frakg=\frakg(\mO)$，$L^{<0}\frakg=t^{-1}\frakg\otimes k[t^{-1}]$。因此 $L^{+}G\times L^{<0}G\to LG$ 是开嵌入。
\end{proof}

现在考虑 $L^-G\cdot t^\mu\cdot \Gr^0$。它是开集，且是 $L^-G$-轨道的并。容易看出 $t^\la$ 属于它，当且仅当 $\Gr^\la\cap \Gr_\mu\neq \emptyset$，故由 (2) 当且仅当 $\la\leq \mu$。因此 $\bigsqcup_{\la\leq \mu}\Gr^\la=L^-G\cdot t^\mu \cdot \Gr^0$ 开。于是 $\Gr^{\geq \mu}$ 闭。还需说明对每个 $\la>\mu$，$\Gr^\la$ 在 $\Gr^\mu$ 的闭包中。如命题 \ref{geom of Sch} 那样，可设 $\la=\mu+\al$，其中 $\al$ 是某个正根。取 \eqref{curve C} 定义的曲线 $C_{\mu+\al,\al}$，则容易看出 $t^{\mu+\al}\in C_{\mu+\al,\al}$，而补集 $\bA^1\subset \Gr^\mu$。证毕。

(4) 由 (3) 与 \eqref{dim opp Sch} 推出。
\end{proof}

\begin{rmk}
(i) 开子集 $\Gr^0$ 通常称为 $\Gr$ 的「大开胞腔」。

(ii) 正如刚才所见，$L^+G$ 与 $L^-G$ 扮演的角色不同，因此考虑商 $LG/L^-G$ 也很自然。由下文定理 \ref{uniformization for P1}，它可解释为 $\bP^1$ 上带一个平凡化的 $G$-主丛的模空间，平凡化取在 $D_0$（$0\in \bP^1$ 处的圆盘）上。文献中有时把 $LG/L^-G$ 称为 $G$ 的\textbf{厚仿射 Grassmann 簇}。注意与 $\Gr$ 不同，它具有整体性质（定义涉及一条整体曲线），并且是（无穷型的）概形。
\end{rmk}

理解对偶 Schubert 胞腔与簇的最好办法是下述定理。以 $\Bun_G(\bP^1)$ 记 $\bP^1$ 上 $G$-丛的模叠，它是 $k$ 上局部有限展示的代数叠（例如见 \cite{So}）。

\begin{thm}\label{uniformization for P1}
设 $G$ 在 $k$ 上连通约化。则有典范同构
\[ [L^-G\backslash \Gr]\simeq \Bun_G(\bP^1).\]
事实上，映射 $\Gr\to \Bun_G(\bP^1)$ 把 $\Gr$ 呈现为 $\Bun_G(\bP^1)$ 上 \'etale 拓扑下的 $L^-G$-主丛。
\end{thm}

\begin{proof}
该论断等价于说：经过 \'etale 覆盖 $R\to R'$ 之后，$\bP^1\times \Spec R$ 上每个 $G$-丛在 $(\bP^1-\{0\})\times \Spec R'$ 上平凡。这由两点推出：(i) 经过可能的 \'etale 基变换 $R\to R'$，$X_R$ 上每个 $G$-丛都可约化到 Borel（\cite[Theorem 1]{DS}）；(ii) 由于 $B$ 有滤过、其伴随分次为若干 $\bG_a$ 与 $\bG_m$，$(\bP^1-\{0\})\times \Spec R$ 上每个 $B$-丛经可能的 \'etale 基变换 $R\to R'$ 后平凡。
\end{proof}

在 $k$-点层面，该定理说
\[ \Bun_{G}(\bP^1)(k)\cong G(k[t^{-1}])\backslash G(k((t)))/ G(k[[t]]).\]
特别地，Grothendieck 关于 $\bP^1$ 上向量丛的定理等价于（$G=\GL_n$ 的）Birkhoff 分解 \eqref{Birk decom}。确切地说，对应于 $t^\mu$ 的双陪集的 $G$-丛，是由线丛 $\mO_{\bP^1}(1)$ 沿 $\bG_m\stackrel{\mu}{\to}T\subset G$ 推出的。把这个 $G$-丛记作 $\mO(\mu)$，并把它在 $\Bun_G(\bP^1)$ 上对应的点记作 $[\mO(\mu)]$。于是命题 \ref{geom of opp Sch} 可改述为

\begin{prop}\label{geom of strata}~
\begin{enumerate}
\item 点 $[\mO(\la)]$ 在 $[\mO(\mu)]$ 的闭包中，当且仅当 $\la\geq \mu$；
\item $[\mO(\mu)]$ 的闭包是维数为 $-(2\rho,\mu)-\dim P_\mu$ 的闭子叠。
\end{enumerate}
\end{prop}

作为命题 \ref{geom of opp Sch} 的应用，还可以引入横截片。

\begin{prop}
$L^{<0}G\cdot t^\la$ 是 $\Gr_\la$ 在 $\Gr$ 中的横截片。更确切地说，对每个 $\la\leq \mu$，令 $S^\la_{\leq \mu}=L^{<0}G\cdot t^\la\cap \Gr_{\leq \mu}$，则 $S^\la_{\leq \mu}$ 与每个 $\Gr_\nu$（$\la\leq \nu\leq \mu$）横截相交。
\end{prop}

\begin{proof}由引理 \ref{big cell} 推出。
\end{proof}

\begin{rmk}\label{symplectic sing}
簇 $S^\la_{\leq \mu}$ 是带辛奇点的簇的例子（见 \cite[Theorem 2.7]{KWWY}），该概念的定义见 \cite[Definition 1.1]{Bea}。由 \cite[Proposition 1.3]{Bea}，辛奇点是有理 Gorenstein 的，这与定理 \ref{singularity} 一致。
\end{rmk}

下面讨论几个例子。

\begin{ex}\label{slice GLn}
考虑 $G=\GL_n$，$\mu=(n,0,\ldots,0)$。首先有

\begin{lem}\label{Lus obs}
映射 \eqref{nil to Gr} 是开嵌入。
\end{lem}

\begin{proof}
事实上，映射 $\mN_n\to L^{<0}\GL_n$，$A\mapsto 1-At^{-1}$，是闭嵌入，故由引理 \ref{big cell}，映射 $\mN_n\to \Gr_{\GL_n}$，$A\mapsto (t-A)\La_0$，是局部闭嵌入且像含于 $\Gr_{\leq \mu}$。因此 \eqref{nil to Gr} 本身是局部闭嵌入。由于 $\Gr_{\leq \mu}$ 不可约且既约、$\dim \mN_n=\dim \Gr_{\leq \mu}=n(n-1)$，\eqref{nil to Gr} 必为开嵌入。
\end{proof}

现取 $\la=(1,\ldots,1)$。注意 $\Gr_\la$ 是一个点。于是引理 \ref{Lus obs} 等价于说 \eqref{nil to Gr} 实现了横截片 $S_{\leq \mu}^\la\subset \Gr_{\leq \mu}$。因此 $S_{\leq \mu}^\nu$ 同构于 $\mO_\nu$ 的 Slodowy 片，其中 $\la\leq \nu\leq \mu$。
\end{ex}

\begin{ex}
设 $G_\der$ 单，$\mu=\theta$ 拟极小。由与上面相同的构造，$S^0_{\leq \theta}$ 同构于 $\frakg$ 的极小幂零轨道闭包 $\overline{\mO}_{\on{min}}$。
\end{ex}

\begin{rmk}
注意定理 \ref{singularity} 蕴含：$\GL_n$ 的幂零轨道以及一般半单李代数 $\frakg$ 的极小幂零轨道都正规、Cohen--Macaulay、Gorenstein，且有有理奇点。这对半单李代数中一般的幂零轨道并不成立，这说明上面横截片与幂零轨道之间的关系不能推广到一般情形。
\end{rmk}
